% ECONOMICS_PROBLEM: {"id": "C73-43", "title": "Approximate stopping equilibrium with nested information", "jel_code": "C73", "source_id": "OP-0166"}
% !TeX program = xelatex
\documentclass[11pt,a4paper]{article}
\usepackage[margin=23mm]{geometry}
\usepackage{fontspec}
\setmainfont{Latin Modern Roman}
\usepackage{amsmath,amssymb,mathtools}
\usepackage{xcolor,enumitem,booktabs,longtable,array,xurl,hyperref}
\definecolor{muted}{HTML}{53636C}
\hypersetup{colorlinks=true,urlcolor=blue,linkcolor=blue}
\setlength{\parindent}{0pt}
\setlength{\parskip}{5pt}
\newcommand{\R}{\mathbb R}
\newcommand{\E}{\mathbb E}
\newcommand{\N}{\mathbb N}
\newcommand{\ind}{\mathbf 1}
\newcommand{\Var}{\operatorname{Var}}
\newcommand{\Cov}{\operatorname{Cov}}
\newcommand{\supp}{\operatorname{supp}}
\DeclareMathOperator*{\argmin}{arg\,min}
\DeclareMathOperator*{\argmax}{arg\,max}
\newcommand{\entrymeta}[3]{{\sffamily\small\color{muted}\textbf{\texttt{#1}}\quad|\quad #2\par #3\par}}
\newcommand{\Problemref}[1]{\texttt{#1}}
\newcommand{\JELref}[2]{\texttt{#2} (JEL \texttt{#1})}
\begin{document}
\section*{Approximate stopping equilibrium with nested information}
\textbf{Problem ID:} \texttt{C73-43}\quad \textbf{JEL:} \texttt{C73}\par
% BEGIN ECONOMICS STATEMENT
\label{problem:OP-0166}
% GAME_THEORY_PROBLEM: {"local_id": "GT-036", "original_number": 36, "kind": "R", "entry_kind": "statement_only", "published": false, "review_required": true, "category": "game_theory"}
\label{game:GT-036}
{\sffamily\small\color{muted}
{\sffamily\small\color{muted}\textbf{GT-036} | Bayesian and stopping games\par R | Fully specified discrete-time representative of a source question\par}
\par}
\subsection*{Definitions and mathematical statement}
Fix a common-prior probability space $(\Omega,\mathcal F,\Pr)$, players $N=\{1,\ldots,n\}$, and dates $t=0,\ldots,T$ with finite $T$. Player $i$ has a filtration $(\mathcal F_{i,t})_{t=0}^{T}$, with
\[
 \mathcal F_{n,t}\subseteq\cdots\subseteq\mathcal F_{1,t}\subseteq\mathcal F
 \qquad(t=0,\ldots,T).
\]
Give each player independent private uniform random seeds independent of $\Omega$. Its strategy is a randomized stopping time $\tau_i\in\{0,\ldots,T,\infty\}$ for its filtration enlarged by its own seeds. The game ends at $\tau=\min_i\tau_i$; if $\tau=t<\infty$, the simultaneous stopper set is $J=\{i:\tau_i=t\}$. For every nonempty $J\subseteq N$ and date $t$, specify bounded measurable payoffs $X_{i,J,t}$; specify bounded $Y_i$ if no player stops. Thus
\[
 U_i(\tau_1,\ldots,\tau_n)
 =\mathbb E\!\left[\sum_{t=0}^{T}\sum_{\varnothing\ne J\subseteq N}
 X_{i,J,t}\mathbf1\{\tau=t,\ \{j:\tau_j=t\}=J\}
 +Y_i\mathbf1\{\tau=\infty\}\right].
\]
Does every game in this specified family satisfy
\[
 \forall\varepsilon>0\quad\exists(\tau_i)_i\quad\forall i\quad\forall\widehat\tau_i,
 \qquad U_i(\widehat\tau_i,\tau_{-i})\leq U_i(\tau)+\varepsilon?
\]
Public survival before stopping is observed; no earlier stopping announcement exists on a surviving history.
\subsection*{Short English statement}
Does nested information alone ensure approximate equilibrium in finite-horizon stopping games?
\subsection*{Variants and scope boundaries}
The probability space need not be finite; finite spaces give a finite Bayesian strategic game. Continuous time requires additional conventions for admissible stopping times, path regularity, and tie payoffs. The stronger condition $\mathcal F_{1,t}\subseteq\mathcal F_{n,t+1}$, used in the source's positive multistage result, is not imposed. Boundedness here is an explicit representative regularity restriction.
\subsection*{Source relationship and status}
[S1], Section 4, poses the stopping-game question under its contemporaneous nesting assumption A3-a. It does not fully specify all discrete- and continuous-time payoff and tie conventions in that discussion. The displayed family resolves those ambiguities editorially; it is not presented as a literal universal theorem statement from the paper.
\subsection*{References}
\begingroup\footnotesize\raggedright
\textbf{[S1]} Royi Jacobovic and Eilon Solan, \emph{Bayesian Games with Nested Information}, arXiv:2402.14450v1 (2024). Locators: Section 4, subsections ``Multi-stage Bayesian games'' and ``Stopping games with asymmetric information,'' including assumptions A3-a and A3-b. \url{https://arxiv.org/html/2402.14450v1}
\par\endgroup

\subsection*{Common conventions: Source mathematical conventions}

\textbf{Finite games.} For a finite set $A$, $\Delta(A)$ is its probability
simplex. Players choose independent mixed strategies in Nash equilibrium;
correlation is allowed only when explicitly part of the solution concept.
In a game with payoff functions $u_i$ and strategy profile $x$, an additive
$\varepsilon$ Nash equilibrium satisfies
\[
 u_i(x)\geq u_i(x_i',x_{-i})-\varepsilon
 \quad\text{for every player $i$ and every admissible deviation $x_i'$}.
\]
For numerical approximation, payoffs are normalized as locally specified.
An $\varepsilon$ payoff residual is distinct from distance at most
$\varepsilon$ to the set of exact equilibria. Point convergence, convergence
of empirical play, and convergence in distance to an equilibrium set are
also distinct.

\textbf{Computational representations.} Unless stated otherwise, a complexity
question takes rational data with binary encoding; polynomial time is measured
in total bit length. A fixed-accuracy polynomial algorithm is not necessarily
polynomial in $1/\varepsilon$, and neither is necessarily polynomial in
$\log(1/\varepsilon)$. Succinct games and value, demand, or sampling oracles
must declare their interfaces. Exact real solutions require an explicit
representation; an unspecified list of arbitrary real numbers is not a
finite computational output. A claimed algorithmic barrier is meaningful
only for its specified model and reductions.

\textbf{Dynamic games.} Histories, information, observation, and allowable
randomization are part of the model. A stationary strategy depends only on
the current state; a positional strategy in a deterministic graph game
chooses one action at each controlled vertex. Nash equilibrium at an initial
state and equilibrium after every history are different requirements.
An expectation of a pathwise $\liminf$ payoff, a $\liminf$ of expected averages,
a limiting expected average, and a uniform finite-horizon equilibrium are
different criteria. None is silently substituted for another.

\textbf{Allocation and mechanisms.} A complete allocation assigns every item
exactly once unless otherwise stated. Zero-valued goods, negative values,
capacity constraints, feasibility, lotteries, and copies of goods affect
fairness definitions and are specified locally. Dominant-strategy incentive
compatibility, Bayesian incentive compatibility, universal truthfulness,
and truthfulness in expectation impose different inequalities. Whenever
an objective ratio has zero denominator, its convention is stated or those
instances are excluded.

\textbf{Infinite-dimensional models.} Expectations, integrals, stochastic
equations, and optimization objectives require their local regularity,
integrability, filtrations, and admissible domains. A backward--forward
mean-field system is a boundary-value problem; it is not an initial-value
semiflow without an additional construction. Long-time, large-population,
and vanishing-noise limits require an order of limits and a convergence
topology. If a minimum need not be attained, the objective is its infimum
or supremum and the approximation requirement is stated separately.
% END ECONOMICS STATEMENT
\end{document}
